\documentclass[11pt]{article}

\usepackage[margin=0.82in]{geometry}
\usepackage{amsmath,amssymb,bm}
\usepackage{microtype}
\usepackage[hidelinks]{hyperref}
\usepackage{enumitem}
\usepackage{xcolor}
\usepackage{booktabs}

\newcommand{\e}{\epsilon}
\newcommand{\ii}{\mathrm{i}}
\newcommand{\bG}{\mathbf{G}}
\newcommand{\bQ}{\mathbf{Q}}
\newcommand{\bK}{\mathbf{K}}
\newcommand{\bk}{\mathbf{k}}
\newcommand{\br}{\mathbf{r}}
\newcommand{\bu}{\mathbf{u}}
\newcommand{\cO}{\mathcal{O}}

\title{\vspace{-1.2cm}Charge-Density-Wave Diffraction from a Commensurate Supercell}
\author{}
\date{}

\title{Challenge 59: Final answer}
\author{}
\begin{document}
\maketitle
Thus
\begin{equation}
    F(\bk)=F_0(\bk)+F_1(\bk)+\cO(\e^2),
\end{equation}
with
\begin{align}
F_0(\bk)
&=f\sum_{j=0}^{M-1}e^{\ii k_xja},\\
F_1(\bk)
&=\ii f(\bk\cdot\bm{\e})
  \sum_{j=0}^{M-1}e^{\ii k_xja}
  \sin\!\left(\frac{2\pi j}{M}\right).
\end{align}

Index the supercell reciprocal points along $x$ by
\begin{equation}
    k_x=\frac{2\pi}{Ma}n_x.
\end{equation}
Then
\begin{equation}
F_0=f\sum_{j=0}^{M-1}e^{\ii2\pi n_xj/M}
=
\begin{cases}
Mf,&n_x=0\pmod M,\\
0,&\text{otherwise}.
\end{cases}
\end{equation}
Hence the ordinary cubic Bragg peaks occur at $n_x=mM$; in the lowest nonzero zone the parent peak is $n_x=M$.


Around the parent peak $n_x=M$, the only first-order CDW satellites are therefore
\begin{equation}
    \boxed{n_x=M-1\quad\text{and}\quad n_x=M+1},
\end{equation}
or equivalently
\begin{equation}
    \boxed{\bk=\bG-\bQ\quad\text{and}\quad\bk=\bG+\bQ}.
\end{equation}

Their structure factors are
\begin{equation}
    \boxed{
    F(\bG-\bQ)
    \simeq
    +\frac{Mf}{2}(\bk\cdot\bm{\e})
    },
\end{equation}
\begin{equation}
    \boxed{
    F(\bG+\bQ)
    \simeq
    -\frac{Mf}{2}(\bk\cdot\bm{\e})
    },
\end{equation}
up to the global Fourier-sign convention. Since $\bm{\e}\parallel\hat{\mathbf{x}}$,
\begin{equation}
    \bk\cdot\bm{\e}=k_x\e.
\end{equation}
Consequently the leading satellite intensity is
\begin{equation}
    \boxed{I_{\pm Q}\propto |F_1|^2\propto(\bk\cdot\bm{\e})^2\propto\e^2.}
\end{equation}

So the $M$-site supercell sum cancels all new reciprocal points to first order except the two roots-of-unity conditions $n_x=M\pm1$, giving the CDW satellites $\bG\pm\bQ$ with amplitude $\pm(Mf/2)(\bk\cdot\bm{\e})$ and intensity proportional to $\e^2$.
\end{document}
